% !TEX root = ../jsr_paper.tex
% Flight-Corpus Integration Test (headline validation section)

The component benchmarks of the preceding sections isolate individual physical
models against canonical wind-tunnel and range data. The integration test asks
the question that matters operationally: when every model is composed inside a
full six-degree-of-freedom trajectory and run end to end against real flights,
how accurately does the predictor recover the measured apogee? This section
reports that result on an externally defined corpus, quantifies its uncertainty
with non-parametric inference, and discloses the in-sample dependence honestly.

\subsection{Corpus Definition and Inclusion Criterion}
\label{sec:corpus-def}

The headline corpus comprises \num{25} flights spanning peak Mach
\numrange{0.54}{4.33}: \num{23} single-stage vehicles together with \num{2}
two-stage flights (the AeroPac~104K flight at Mach~3.04 and the MESOS~293K
flight at Mach~4.33), drawn directly from the public RASAero~II
flight-comparison set assembled by Rogers~\cite{rogers2015,
rogers_rasaero_alt}. The MESOS~293K flight closes the corpus at its highest
Mach. The inclusion criterion is the central guard against
selection bias: every flight is taken from a comparison set chosen and published
by an \emph{external} party for an independent purpose, and the set is admitted
in its entirety. No flight was added, removed, or reweighted on the basis of how
well the present predictor happened to perform on it. The accuracy statistics
below are therefore an outcome-independent validation result rather than a
curated demonstration. Vehicles range from \SI{1.75}{in} hobby airframes to the
\SI{8.0}{in} Qu8k and the \SI{6.13}{in} A-601 Kinsel, and ground-truth apogees
derive from heterogeneous instrumentation (barometric altimetry, optical
tracking, integrated accelerometry, and GPS), as recorded in
Table~\ref{tab:corpus-perflight}.

A separate set of high-Mach sounding rockets reaching Mach~7 is treated as an
\emph{exploratory} study in Sec.~\ref{sec:hypersonic}, deliberately segregated
from this headline validation; the present corpus is supersonic-validated and
contains no hypersonic flight (maximum Mach~4.33).

\subsection{Headline Accuracy}
\label{sec:corpus-headline}

Across the \num{25} flights, the OpenRocket-Plus apogee-prediction error (ORP,
defined as $100\,(h_\text{pred}-h_\text{meas})/h_\text{meas}$) has a mean signed
value of \SI{-0.38}{\percent} with sample standard deviation
$\sigma=\SI{5.44}{\percent}$, root-mean-square error of \SI{5.34}{\percent}, and
mean absolute error of \SI{4.74}{\percent}. All \num{25} of \num{25} flights
fall within \SI{\pm 10}{\percent}, and \num{14} of \num{25} within
\SI{\pm 5}{\percent}; the largest single residual is \SI{8.7}{\percent}.
Table~\ref{tab:corpus-headline} summarizes these figures alongside the recorded
RASAero~II predictions for the same flights.

Because the corpus deviates from normality (Sec.~\ref{sec:corpus-dist}), every
interval below is a 95\% non-parametric percentile bootstrap confidence
interval (\num{20000} resamples, fixed seed \texttt{0x51A7EA}). The mean-error
interval is $[-2.41,+1.72]\,\si{\percent}$; it brackets zero, so no bias is
detectable on this corpus. We are careful not to overstate this as proven
unbiasedness: at $n=25$ the interval is roughly \SI{\pm 2}{pp} wide, so a small
true bias would not be resolved, and bracketing zero is an absence of evidence
for bias rather than evidence of its absence. The remaining intervals are
$[4.20,6.25]\,\si{\percent}$ for $\sigma$, $[4.38,6.19]\,\si{\percent}$ for RMSE,
and $[3.78,5.70]\,\si{\percent}$ for MAE. The within-\SI{\pm 5}{\percent} rate of
\SI{56}{\percent} carries an interval of $[36,76]\,\si{\percent}$. The
within-\SI{\pm 10}{\percent} rate is \SI{100}{\percent} (\num{25} of
\num{25}); because the bootstrap is degenerate for an all-pass proportion
(every resample contains only passing flights, returning $[100,100]$ as an
artifact), we instead report an exact binomial Clopper--Pearson interval,
whose lower bound is \SI{86.3}{\percent} at \SI{95}{\percent} two-sided
confidence. Decomposing the mean
squared error into squared bias and variance yields a $\text{bias}^2/\text{MSE}$
ratio of \num{0.01} for ORP, against \num{0.16} for RASAero~II: the ORP residual
is essentially pure variance. Operationally this means the remaining error is
dominated by irreducible flight-to-flight scatter---build tolerance, motor lot
variation, and atmosphere---rather than a directional defect in the
aerodynamic model. Figure~\ref{fig:corpus-bias-variance} visualizes this
decomposition and Fig.~\ref{fig:corpus-error-mach} shows the residuals against
peak Mach.

\begin{table}[t]
  \centering
  \caption{Headline accuracy on the externally defined \num{25}-flight corpus
    (Mach \numrange{0.54}{4.33}). RASAero~II values are Rogers' recorded
    predictions for the same flights.}
  \label{tab:corpus-headline}
  \begin{tabular}{lccccccc}
    \toprule
    Predictor & $N$ & Mean & $\sigma$ & RMSE & MAE & $\le\!\pm5\%$ & $\le\!\pm10\%$ \\
    \midrule
    OpenRocket-Plus & 25 & $-0.38$ & 5.44 & 5.34 & 4.74 & 14/25 & 25/25 \\
    RASAero~II      & 25 & $+2.46$ & 5.81 & 6.20 & 5.34 & 13/25 & 22/25 \\
    \bottomrule
  \end{tabular}
  \\[2pt]
  {\footnotesize All entries in percent except $N$ and the count columns.}
\end{table}

\begin{figure}[t]
  \centering
  \includegraphics[width=\columnwidth]{error_vs_mach}
  \caption{Apogee-prediction error versus peak Mach for the \num{25}-flight
    corpus. Residuals remain bounded within \SI{\pm 10}{\percent} across the
    full subsonic-to-supersonic range; the transonic band carries the largest
    negative bias.}
  \label{fig:corpus-error-mach}
\end{figure}

\begin{figure}[t]
  \centering
  \includegraphics[width=\columnwidth]{bias_variance}
  \caption{Bias--variance decomposition of the mean squared apogee error.
    For OpenRocket-Plus the squared-bias term is a negligible fraction
    ($\text{bias}^2/\text{MSE}=0.01$) of the total, indicating an
    essentially unbiased predictor whose residual is dominated by variance.}
  \label{fig:corpus-bias-variance}
\end{figure}

\subsection{Per-Regime Behavior}
\label{sec:corpus-regime}

Partitioning the corpus by Mach regime (Table~\ref{tab:corpus-regime}) localizes
the error structure. The subsonic ($M<0.8$, $n=9$) and low-supersonic
($1.3\le M<3.0$, $n=5$) bands carry small positive biases (\SI{+2.54}{\percent}
and \SI{+1.83}{\percent}). The transonic band ($0.8\le M<1.3$, $n=7$) is the
disclosed regime weakness, with a mean bias of \SI{-3.66}{\percent}: this is the
Mach range where shock-induced drag rise is most sensitive to geometry detail
and where the present models are weakest. The high-supersonic band
($3.0\le M<5.0$, $n=4$) has a \SI{-3.98}{\percent} bias but the tightest scatter
($\sigma=\SI{2.97}{\percent}$). Above Mach~3 the sample is small ($n\le 4$ per
band), so these figures are reported descriptively; no inferential claim is made
there. The hypersonic regime ($M>5.0$) is empty by construction. The regime
breakdown is shown graphically in Fig.~\ref{fig:corpus-regime}.

\begin{table}[t]
  \centering
  \caption{Per-regime ORP accuracy on the \num{25}-flight corpus. Biases,
    $\sigma$, RMSE, and MAE in percent; count columns are flights within
    tolerance.}
  \label{tab:corpus-regime}
  \begin{tabular}{lccccccc}
    \toprule
    Regime & Mach & $N$ & Bias & $\sigma$ & RMSE & $\le\!5\%$ & $\le\!10\%$ \\
    \midrule
    Subsonic        & $<0.8$        & 9 & $+2.54$ & 4.37 & 4.83 & 7 & 9 \\
    Transonic       & $0.8$--$1.3$  & 7 & $-3.66$ & 5.36 & 6.17 & 2 & 7 \\
    Low supersonic  & $1.3$--$3.0$  & 5 & $+1.83$ & 5.70 & 5.42 & 3 & 5 \\
    High supersonic & $3.0$--$5.0$  & 4 & $-3.98$ & 2.97 & 4.73 & 2 & 4 \\
    Hypersonic      & $>5.0$        & 0 & ---     & ---  & ---  & --- & --- \\
    \bottomrule
  \end{tabular}
\end{table}

\begin{figure}[t]
  \centering
  \includegraphics[width=\columnwidth]{regime_breakdown}
  \caption{Per-regime error distribution. The transonic band carries the
    largest negative bias; the high-supersonic band has the tightest scatter.
    Bands above Mach~3 contain four or fewer flights and are descriptive only.}
  \label{fig:corpus-regime}
\end{figure}

\subsection{Distribution and Choice of Inference}
\label{sec:corpus-dist}

The choice of non-parametric inference is dictated by the residual
distribution, not adopted by convenience. The Shapiro--Wilk test rejects
normality of the \num{25} ORP residuals ($W=0.905$, $p=0.023$ at
$\alpha=0.05$), and the Anderson--Darling statistic ($A^2=0.922$) likewise
exceeds its critical value ($0.728$). The residuals are mildly right-skewed
(skewness $+0.43$) and distinctly platykurtic (excess kurtosis $-1.14$): the
distribution is light-tailed rather than heavy-tailed, consistent with the
bounded maximum residual of \SI{8.7}{\percent}. We therefore adopt the Wilcoxon
signed-rank test as the primary hypothesis test and bootstrap confidence
intervals in place of normal-theory intervals throughout. We also caution that,
at $n=25$, normality tests have limited power; failing to reject normality would
not establish it, which is a further reason to prefer distribution-free methods.
Figure~\ref{fig:corpus-distributions} compares the ORP and RASAero~II error
distributions, and Fig.~\ref{fig:corpus-bootstrap} shows the bootstrap sampling
distributions underlying the headline intervals.

\begin{figure}[t]
  \centering
  \includegraphics[width=\columnwidth]{predictor_distributions}
  \caption{Error distributions for OpenRocket-Plus and RASAero~II on the
    \num{25}-flight corpus. Both are light-tailed; the ORP distribution centers
    near zero.}
  \label{fig:corpus-distributions}
\end{figure}

\begin{figure}[t]
  \centering
  \includegraphics[width=\columnwidth]{bootstrap_distributions}
  \caption{Bootstrap sampling distributions (\num{20000} resamples, seed
    \texttt{0x51A7EA}) for the headline ORP statistics. The mean-error
    distribution straddles zero; no bias is detectable at this sample size.}
  \label{fig:corpus-bootstrap}
\end{figure}

\subsection{Paired Comparison Against RASAero~II}
\label{sec:corpus-paired}

Because each flight carries a recorded RASAero~II prediction, the two methods
can be compared head to head on identical vehicles and motors. On absolute
error, the paired mean difference $|\text{ORP}|-|\text{RAS}|$ is
\SI{-0.60}{pp} with a 95\% non-parametric percentile bootstrap confidence
interval of $[-2.16,+0.96]\,\si{pp}$, which straddles zero. The Wilcoxon
signed-rank test on the paired absolute errors returns $W=143.0$, $p=0.615$:
no significant difference in accuracy is detected on this corpus, with ORP closer
on \num{13} flights and RASAero~II closer on \num{12}. We read this as a failure
to reject rather than as positive proof of equivalence: at $n=25$ the paired test
has limited power, so $p=0.615$ alone is absence of evidence for a difference, not
evidence of its absence. Equivalence is supported more defensibly by the bootstrap
interval on the paired mean above, which lies entirely within roughly
\SI{\pm 2}{pp}---a two-one-sided-test argument against any difference larger than
that margin---than by the non-significant Wilcoxon $p$-value. The Bland--Altman limits
of agreement are \SI{\pm 14.3}{\percent} about a mean offset of
\SI{-2.84}{\percent}. The honest claim supported by these statistics is
\emph{parity}, not superiority. We further disclose an important asymmetry: the
RASAero~II values are Rogers' published recorded predictions for these flights,
not fresh pre-flight reruns we generated and can independently verify, so the
comparison is to a version-locked external record rather than a controlled
re-execution. In addition, for the two two-stage flights---AeroPac~104K and
MESOS~293K---the recorded RASAero~II apogees are \emph{post-flight} simulations
in which the ignition delay and launch angle were adjusted to match the measured
flight; the head-to-head on these two flights is therefore not a blind forward
prediction, and the corresponding RASAero~II errors should be read accordingly.
Figure~\ref{fig:corpus-paired} shows the paired residuals.

\begin{figure}[t]
  \centering
  \includegraphics[width=\columnwidth]{predictor_paired}
  \caption{Paired per-flight comparison of OpenRocket-Plus and RASAero~II
    apogee errors. The paired absolute-error difference is statistically
    indistinguishable from zero (Wilcoxon $p=0.615$), supporting a parity
    conclusion.}
  \label{fig:corpus-paired}
\end{figure}

\subsection{In-Sample Dependence and Decontaminated Holdout}
\label{sec:corpus-insample}

Two base-drag scale constants are calibrated at the corpus level and are
therefore not free of in-sample dependence: \texttt{THICK\_BL\_K}~$=2.2$
(anchored to the Raven flight) and \texttt{SLENDER\_BODY\_K}~$=0.0025$ (anchored
to Raven, Rabia, and Kinsel, with Torrent used as a source diagnostic). The
headline statistics are thus partly in-sample, and we report this rather than
obscure it.

The primary defense against overfitting is a decontaminated prospective holdout.
Every flight that either of the two constants touched---Raven, Rabia, Rabia Short
Fin Can, Kinsel, and Torrent---is placed in the development partition, leaving a
genuinely blind holdout. On the development partition ($n=13$) the mean error is
\SI{+0.22}{\percent} with MAE \SI{5.47}{\percent}; on the blind holdout
($n=12$) the mean error is \SI{-1.03}{\percent} with MAE \SI{3.95}{\percent}.
The holdout is \emph{more} accurate than the development set, which is the
signature of constants that generalize rather than overfit. (An earlier split
left Rabia in the holdout, contaminating it; the decontaminated split corrects
this.) This prospective generalization, not any operational sweep, is the basis
on which we rest the in-sample disclosure.

\subsection{The Two-Stage High-Mach Closure (MESOS 293K)}
\label{sec:corpus-mesos}

The MESOS~293K flight closes the corpus at Mach~4.33 and is one of the two
two-stage flights in the corpus (alongside AeroPac~104K), exercising staging and
coast-phase modeling not present in the single-stage flights. It also carries the
single largest per-flight error in the corpus: the current archived code release
predicts an apogee of \SI{273056}{ft}, an error of \SI{-6.96}{\percent} against
the measured \SI{293488}{ft}. This value is reproducible from the archived code
release---verified by running the MESOS flight in isolation---and is the value
reported throughout this paper and carried in the published database. As the only
high-Mach two-stage closure, MESOS~293K simultaneously exercises stage
separation, coast-phase aerodynamics, and Mach-4 base drag, and so bounds the
framework's accuracy at the upper edge of its validated envelope; we do not
isolate which of these compounding effects dominates the error. The flight
remains within the \SI{\pm 10}{\percent} envelope and is retained in the corpus
under the same outcome-independent inclusion rule applied to every other flight;
we report the \SI{-6.96}{\percent} value without adjustment and carry it forward
to the limitations discussion.

\subsection{Per-Flight Results}
\label{sec:corpus-perflight}

Table~\ref{tab:corpus-perflight} lists all \num{25} flights with vehicle, motor,
peak Mach, measured and predicted apogees, and per-predictor error, ordered by
peak Mach.

\begin{table*}[t]
  \centering
  \caption{Per-flight results for the \num{25}-flight corpus, ordered by peak
    Mach. Apogees in feet; errors in percent of measured apogee. RASAero~II
    values are Rogers' recorded predictions for the same flights.}
  \label{tab:corpus-perflight}
  \begin{tabular}{rlcrrrrrr}
    \toprule
    & Vehicle & Motor & $M_\text{pk}$
      & $h_\text{meas}$ & $h_\text{RAS}$ & $h_\text{ORP}$
      & $\varepsilon_\text{RAS}$ & $\varepsilon_\text{ORP}$ \\
    \midrule
    1  & Thunder \& Lightning            & I284W           & 0.54 & 3577   & 3989   & 3877   & $+11.52$ & $+8.4$  \\
    2  & Gibb                            & I284W           & 0.55 & 3913   & 4310   & 3989   & $+10.15$ & $+1.9$  \\
    3  & Cancer Descending               & M1297W          & 0.56 & 6188   & 6328   & 6044   & $+2.26$  & $-2.3$  \\
    4  & EZI-65 J450ST                   & J450ST          & 0.60 & 3965   & 4214   & 4158   & $+6.28$  & $+4.9$  \\
    5  & Caliber Isp 04 AVTC Team 2      & I205            & 0.64 & 3710   & 3871   & 3890   & $+4.34$  & $+4.9$  \\
    6  & Caliber Isp 04 AVTC Team 3      & I205            & 0.64 & 3964   & 3871   & 3889   & $-2.35$  & $-1.9$  \\
    7  & Caliber Isp 04 AVTC Team 1      & I205            & 0.66 & 3837   & 3943   & 3960   & $+2.76$  & $+3.2$  \\
    8  & Byrum                           & J570W           & 0.75 & 5732   & 5280   & 6161   & $-7.89$  & $+7.5$  \\
    9  & Ion Drive                       & K550W           & 0.79 & 8027   & 8642   & 7730   & $+7.66$  & $-3.7$  \\
    10 & Caliber Isp 05 ARO-414 (Disc.)  & I285            & 0.81 & 4930   & 4831   & 4772   & $-2.01$  & $-3.2$  \\
    11 & Blister                         & K1075GG         & 0.83 & 9026   & 8347   & 8268   & $-7.52$  & $-8.4$  \\
    12 & Caliber Isp 05 ARO-414 (Col.)   & I285            & 0.84 & 5085   & 4842   & 4777   & $-4.78$  & $-6.1$  \\
    13 & Rabia -- Short Fin Can          & L730            & 0.86 & 10584  & 10376  & 9916   & $-1.97$  & $-6.3$  \\
    14 & Raven                           & J570W           & 1.07 & 8815   & 9288   & 9489   & $+5.37$  & $+7.6$  \\
    15 & Rabia                           & L1080BB         & 1.14 & 12745  & 12777  & 11913  & $+0.25$  & $-6.5$  \\
    16 & Torrent                         & M1850GG         & 1.22 & 12807  & 13852  & 12455  & $+8.16$  & $-2.8$  \\
    17 & Kline-Rogers L500               & L500            & 1.98 & 24771  & 26485  & 24179  & $+6.92$  & $-2.4$  \\
    18 & A-601 Kinsel                    & P4935           & 2.19 & 42771  & 41086  & 46499  & $-3.94$  & $+8.7$  \\
    19 & Full Metal Jacket -- BALLS 005  & O10000          & 2.31 & 37981  & 38820  & 37256  & $+2.21$  & $-1.9$  \\
    20 & Full Metal Jacket -- Black Rock 6 & O10000        & 2.46 & 30038  & 32646  & 29239  & $+8.68$  & $-2.7$  \\
    21 & Proteus 6                       & P9381           & 2.87 & 85067  & 86799  & 91339  & $+2.04$  & $+7.4$  \\
    22 & AeroPac 104K Two-Stage          & N1048 / M685W   & 3.04 & 104659 & 113786 & 103602 & $+8.72$\textsuperscript{\dag} & $-1.0$  \\
    23 & Don't Debate This               & N5800           & 3.04 & 56573  & 62308  & 53150  & $+10.14$ & $-6.1$  \\
    24 & Qu8k                            & Q18000          & 3.46 & 121478 & 116254 & 119187 & $-4.30$  & $-1.9$  \\
    25 & MESOS 293K                      & O4374 / M787    & 4.33 & 293488 & 289789 & 273056 & $-1.26$\textsuperscript{\dag} & $-6.96$ \\
    \bottomrule
  \end{tabular}
  \\[2pt]
  {\footnotesize \textsuperscript{\dag}For the two-stage flights AeroPac~104K
    and MESOS~293K, the recorded RASAero~II apogees are post-flight simulations
    in which the ignition delay and launch angle were adjusted to match the
    measured flight; these RASAero~II errors are therefore not blind forward
    predictions.}
\end{table*}
